% !TEX root = hw14-answers.tex
\chead{\textbf{Exercises week 14}}
\begin{document}
\flushleft\textbf{Topic: Causal Discovery}

\begin{enumerate}
\item (2 points) Two key assumptions of Randomized Controlled Trials with treatment variable $C$ and outcome variable $X$ are:
\begin{enumerate}
	\item outcome $X$ does not cause treatment $C$;
	\item outcome $X$ and treatment $C$ are not confounded, i.e., the values for
	 the treatment variable are assigned independently of other (latent) factors that may influence
	 the outcome.
  \end{enumerate}

Think like a conspiracy theorist and describe a situation in which the first assumption is not valid. Do the same for the second assumption.

\begin{gradedsolution}
	For example:
	\begin{enumerate}
		\item Consider the system $X_1 \to C \to X_2$, where $X_1, X_2$ are measurements of health. Based on the measurement $X_1$ the patient is assigned treatment $C$. After a period of time, the outcome $X_2$ is measured. However, the doctors overseeing this RCT mistakenly mix up the columns $X_1$ and $X_2$, and report $X_1$ as outcome.\\
		Another example would be: the patient obtains treatment $C$ after which we measure the unfavourable outcome $X$. A time traveller decides to take action, and changes the value of $C$ (based on $X$) to improve the outcome.
		\item If a pharmaceutical company wants to have a favourable outcome for a new medicine, they can setup a fraudulent RCT where they only provide patients with a high probability of recovery (patients that are mildly ill) with treatment. In this way, the measured illness is a cause of both treatment and outcome.
	\end{enumerate}
\end{gradedsolution}

\item (3 points) Prove Proposition \ref*{prp:y-structures} on Y-structures by using the global Markov property and the faithfulness assumption, and showing that the only (cyclic or acyclic) graphs that are compatible with the observed
conditional independences are the ones in Figure \ref*{fig:Ystruct}, similarly as how is done in the proof of LCD (Proposition \ref*{prp:LCD}).

\begin{gradedsolution}
	We have the conditional (in)dependencies:

	$$\begin{array}{lll}
		X_1 \nIndep X_4,  & X_2 \nIndep X_4, & X_1 \Indep X_2, \\
		X_1 \Indep X_4 \given X_3, & X_2 \Indep X_4 \given X_3, & X_1 \nIndep X_2 \given X_3,
	  \end{array}$$
	
	By faithfulness assumption, from $X_1 \Indep X_2$, $X_1 \Indep X_4 \given X_3$ and $ X_2 \Indep X_4 \given X_3$ it respectively follows that there can be no direct edge between $X_1$ and $X_2$, between $X_1$ and $X_4$ and between $X_2$ and $X_4$, because otherwise the corresponding $\sigma$-separations do not hold. However, since $X_1 \nIndep X_4$ and  $X_2 \nIndep X_4$, there must exists a walk from $X_1$ to $X_4$ and a walk from $X_2$ to $X_4$, which by necessity must go through the only other node $X_3$. These facts together imply that there can only be edges between $X_3$ and the other nodes, and the graph is connected.

	Since $X_1 \Indep X_2$ and $X_1 \nIndep X_2 \given X_3$, it follows that $X_3$ is a collider on the walk from $X_1$ to $X_2$. Therefore, any edges between $X_1$ and $X_3$ and $X_2$ and $X_3$ must have an arrowhead pointing into $X_3$.

	Finally, since $X_1 \nIndep X_4$ and $X_1 \Indep X_4 \given X_3$, conditioning on $X_3$ must block the path from $X_1$ to $X_4$, therefore it cannot be a collider on this walk. As we have shown, edges between $X_1$ and $X_3$ must have an arrow pointing into $X_3$. Therefore, the edge between $X_4$ and $X_3$ cannot have an arrow pointing into $X_3$. It follows that $X_3 \to X_4$.
	
	With these facts taken together, it is easy to see that the only possible graphs are those drawn in Figure \ref*{fig:Ystruct}.
\end{gradedsolution}

\end{enumerate}
\end{document}
